\section{Related Work}
\label{sec:related}
This section relates AgentWorkShop to plant-facing industrial agents and to
production and automation abstractions for context, tool access, and
operational constraints. Table~\ref{tab:nearest} compares integration and
execution features described in selected industrial-agent architectures,
including systems with adjacent but distinct process-control objectives. The
comparison is focused rather than exhaustive; a feature marked as not
specified was not described in the cited study and is not inferred to be
absent.

\subsection{Industrial Agents and Plant Integration}
Industrial-agent systems, including large language model (LLM)-based systems,
span end-to-end automation control, supervisory
services, programmable logic controller (PLC) interfaces, manufacturing
coordination, and operator-facing machine access. Xia et al. present an
LLM agent framework for end-to-end industrial automation control, using
structured prompts and event-driven context for production planning and
operation \cite{xia2025control}. Huang et al. present a multimodal
architecture whose stated focus is automated control generation for
industrial cyber-physical systems \cite{huang2026made}. Biskupovic et al.
study agent-based supervision for service-oriented industrial systems
\cite{biskupovic2025supervision};
Agents4PLC and InstructMPC address PLC code generation and human--LLM
context-aware control, respectively
\cite{liu2026agents4plc,wu2025instructmpc}. Manufacturing studies consider
dynamic resource-capability exploration \cite{lim2025llmenabled}, multi-agent
shop-floor coordination \cite{zhao2026shopfloor}, and digital-twin-supported
fault handling \cite{gill2025llmdt}. Natural-language machine access
\cite{hofmann2025nlc} and real-time LLM control of a robot-operated serial
line \cite{waseem2026pretrained} extend this work to operator interfaces and
robotic-line control. Together, these studies address different operational
contexts, interfaces, and allocations of responsibility between agents,
operators, and controllers.

The closest architecture-level comparison is Xia et al., who connect cognitive,
digital-twin, and physical-automation layers and organize twin information as
state snapshots, event messages, and plan sequences; their case studies
address robotic automation and process simulation
\cite{xia2025architecture}. Other nearby designs expose factory-module
operations through MCP with state-based tool filtering, plan constrained
manufacturing steps, or compose semantic recipes from OPC UA skills
\cite{koehle2026llmdriven,silva2025mcpmanufacturing,koehle2025semanticrecipes}.
These studies describe important parts of plant-facing integration. The
recent systems in Table~\ref{tab:nearest} also combine agent execution with
approval, industrial data, or audit records. Graul et al. ground proposals in
the Asset Administration Shell and route approved actions through OPC UA
Methods, with signed approvals and audit logging
\cite{graul2026aas}. Schall separates advisory language-model agents from
P\&ID-grounded verification and execution, and evaluates structural checks
across five simulated process scenarios
\cite{schall2026safe}. Lin and Chen build a governed sensor--MES--ERP digital
thread for closed-loop monitoring and control, evaluated on real manufacturing
records \cite{lin2026agentic}. These works address asset-specific approval,
process validation, or closed-loop analytics. AgentWorkShop's narrower
architectural delta is the composition of task/chat execution and direct
operation calls with registered harnesses, production context, semantic
drivers, shared per-call dispatch, and service-owned operation records. The
comparison does not imply that other systems lack these capabilities. R20 and
R21 test selected routes through this composition; neither tests comparative
utility or generality beyond the configured paths.
\begin{table*}[!t]
\centering
\caption{Feature-level comparison of selected plant-facing integration
studies. ``Not specified'' means that the cited study does not describe the
feature; it is not evidence of absence. The table compares stated designs and
evaluation evidence, not performance.}
\label{tab:nearest}
\small
\setlength{\tabcolsep}{3pt}
\begin{tabularx}{\textwidth}{@{}p{1.8cm}p{3.0cm}p{2.4cm}p{2.8cm}X@{}}
\toprule
System & Task and production context & Per-call operation checks & Device-facing route & Reported evaluation and record link \\
\midrule
Xia et al. \cite{xia2025architecture} & Cognitive, twin, and physical layers; snapshots, events, and plans & Not specified & Digital-twin/physical-automation integration & Robotic automation and process-simulation case studies \\
K{\"o}hle et al. \cite{koehle2026llmdriven} & Factory-module agents with injected live state & State-based tool filtering & Per-module MCP exposes OPC UA methods; MQTT coordination & Nine simulated production challenges \\
Vieira da Silva et al. \cite{silva2025mcpmanufacturing} & LLM plans manufacturing steps under constraints & Not specified & MCP-exposed manufacturing resource functions & Laboratory multi-step process with constraint handling \\
K{\"o}hle et al. \cite{koehle2025semanticrecipes} & AI-assisted recipes composed from OPC UA skills & Not specified & OPC UA skills and virtual commissioning & Unity assembly-workstation proof of concept \\
Graul et al. \cite{graul2026aas} & AAS-grounded asset and production context & Signed HITL approval; immutable audit log & PLC-validated OPC UA Methods & Two laboratory PoCs; single scripted runs \\
Schall \cite{schall2026safe} & P\&ID, SCADA, time-series, and SOP context & Deterministic topology and rule checks & Advisory proposal, gated execution & Five simulated scenarios; 43 injected proposals and 50 LLM runs \\
Lin and Chen \cite{lin2026agentic} & Sensor--MES--ERP digital thread & SOP-, risk-, and capability-aware gating & Edge monitoring and closed-loop actions & 25,275 real records; monitoring and operational metrics \\
\midrule
This work & Collaborative tasks and direct calls; per-turn line--product--recipe--run context & Scope, policy, grant/binding at dispatch; route-specific state rechecked before write & Semantic node resolves to configured DAQ/DCW driver & R20: 6/6 live-agent traces; R21: 2/2 deterministic fixture traces across two mappings; no process attainment claim \\
\bottomrule
\end{tabularx}
\end{table*}

The evaluation unit is an execution record whose required identifiers link
task, approval, active run, verified receipt, and subsequent sample; other
service records may expose only the identifiers appropriate to their route.
This record path is a concrete implementation focus, not a claim that prior
systems cannot provide equivalent linkage. AgentWorkShop adds no new process
control or safety-verification method. Authority checks and command
confirmation complement safety-aware decision support
\cite{jiang2026agentic}; they do not replace process protection or establish
operator effectiveness.

\subsection{Production Context, Agent Interfaces, and Automation}
Industrial integration also depends on abstractions that are not specific to
agent architectures. MCP standardizes tool discovery and calls with
request-scoped authorization \cite{mcp2024}, whereas
ISA-88 and ISA-95 describe recipes and production operations
\cite{isa88_2010,isa95_2025}. Run-to-run methods connect process observations
to successive production runs \cite{moyne2000r2r}. At the agent level,
AgentSpec and CaMeL address tool behavior and capability-based information
flow \cite{wang2025agentspec,debenedetti2025camel}. Together, these works
provide complementary abstractions for production context, tool interaction,
and information flow. AgentWorkShop composes such abstractions with its task
runtime, operation checks, semantic I/O, and route-specific records; this
paper evaluates selected paths through that composition rather than proposing
a replacement for the underlying standards.

Industrial security and record requirements further constrain that
architecture: FDA guidance discusses trustworthy records and audit trails
\cite{part11}, IEC~62443-3-3 specifies system-security requirements
\cite{iec62443-3-3}, and DNP3 includes select-before-operate
\cite{ieee1815dnp3}. Control-assurance methods address state evolution through
shielding, control-barrier functions, or Simplex architectures
\cite{alshiekh2018shielding,ames2017cbf,sha2001simplex}. These standards and
methods inform the boundaries within which an industrial-agent framework must
operate; they are not contributions claimed by AgentWorkShop.

The evaluation tests selected software paths and simulator scenes, not
physical-plant performance.
